%% LunarAtlas: Draft 5 (Complete Revised Manuscript)
%% Elsevier two-column style (cas-dc class)
%% Authors: Lovekesh Anand, Dua Saeed
%% April 2026

\documentclass[a4paper,twocolumn,10pt]{article}
\usepackage[margin=1.9cm,top=2.6cm,bottom=2.6cm,columnsep=0.6cm]{geometry}
\usepackage{authblk}

% -----------------------------------------------------------------------
%  Core packages
% -----------------------------------------------------------------------
\usepackage{amsmath,amssymb}
\usepackage{siunitx}
\usepackage{graphicx}
\usepackage{booktabs}
\usepackage{tabularx}
\usepackage{multirow}
\usepackage{hyperref}
\usepackage{url}
\usepackage{enumitem}
% caption loaded below with options
\usepackage{subcaption}
\usepackage{float}
\usepackage{xcolor}
\usepackage{microtype}
\usepackage{lineno}   % line numbers for review
\usepackage{stfloats} % allows figure* at bottom of page
\usepackage[font=small,labelfont=bf,justification=centering]{caption}

% -----------------------------------------------------------------------
%  SI units configuration
% -----------------------------------------------------------------------
\sisetup{
  detect-all,
  group-separator={,},
  separate-uncertainty=true
}

% -----------------------------------------------------------------------
%  Hyperref colours (Elsevier-like)
% -----------------------------------------------------------------------
\hypersetup{
  colorlinks=true,
  linkcolor=blue!60!black,
  citecolor=blue!60!black,
  urlcolor=blue!60!black
}

% -----------------------------------------------------------------------
%  Elsevier-like title/abstract styling
% -----------------------------------------------------------------------
\usepackage{titlesec}
\usepackage{fancyhdr}
\usepackage{abstract}
\usepackage{etoolbox}

% Section formatting — Elsevier style
\titleformat{\section}{\normalfont\bfseries\fontsize{10}{12}\selectfont}{\thesection.}{0.5em}{}
\titleformat{\subsection}{\normalfont\bfseries\fontsize{9.5}{11}\selectfont}{\thesubsection}{0.5em}{}
\titleformat{\subsubsection}{\normalfont\itshape\fontsize{9}{11}\selectfont}{\thesubsubsection}{0.5em}{}
\titlespacing*{\section}{0pt}{10pt plus 2pt}{4pt plus 1pt}
\titlespacing*{\subsection}{0pt}{8pt plus 1pt}{3pt plus 1pt}
\titlespacing*{\subsubsection}{0pt}{6pt}{2pt}

% Header/footer
\pagestyle{fancy}
\fancyhf{}
\fancyhead[L]{\small\textit{L.~Anand, D.~Saeed / LunarAtlas: Chandrayaan-3 LIBS Infrastructure}}
\fancyhead[R]{\small\textit{April 2026}}
\fancyfoot[C]{\small\thepage}
\renewcommand{\headrulewidth}{0.3pt}

% Abstract styling
\renewcommand{\abstractnamefont}{\normalfont\bfseries\small}
\renewcommand{\abstracttextfont}{\normalfont\small}
\setlength{\absleftindent}{0pt}
\setlength{\absrightindent}{0pt}

% Paragraph spacing
\setlength{\parskip}{3pt}
\setlength{\parindent}{12pt}

% =========================================================================
%  Document begins
% =========================================================================
\begin{document}

\twocolumn[{%
\begin{@twocolumnfalse}


{\centering
{\LARGE\bfseries LunarAtlas: From PDS4 Archives to\\[4pt]
Analysis-Ready Spectra --- A Reproducible\\[4pt]
Infrastructure for Chandrayaan-3 LIBS Level-1 Data}\par
\vspace{8pt}
{\normalsize\bfseries Lovekesh Anand\textsuperscript{a},\quad Dua Saeed\textsuperscript{a}}\\[3pt]
{\small\textit{\textsuperscript{a}Independent Researchers, India}}
\par
}

\vspace{10pt}
\begin{abstract}
\noindent
We present \textbf{LunarAtlas}, a reproducible data-processing
infrastructure that transforms publicly available Chandrayaan-3
Laser-Induced Breakdown Spectroscopy (LIBS) Level-1 (L1) products
into cleaned, machine-accessible, long-form spectral records suitable
for quantitative lunar science.
Starting from PDS4-compliant L1 archives released by ISRO, the system comprises three core components:
(i)~\textbf{preprocessing} (PDS4 parsing, wide-to-long reshaping, background subtraction with MD5-tracked provenance);
(ii)~\textbf{peak/element identification} (SNR filtering and Gaussian fitting against a curated NIST line list)---presented as a proof-of-concept on a single archival file; and
(iii)~\textbf{adaptive visualization} (min--max downsampling with a zoom-saturation criterion that preserves narrow emission-line features).
Applied to \texttt{ch3\_lib\_002\_20230825T104221\_00\_l1.csv}, the pipeline yields four cleaned measurements
(\textbf{Measurement~IDs~1--4}) with a baseline suppression factor of $\sim$7.1.
Spectral features consistent with Mg, Ca, Al, Fe, and Si emission are detected at high SNR;
these are discussed as illustrative identifications pending broader multi-file validation.
\end{abstract}
\vspace{4pt}
\noindent\textbf{Keywords:} Laser-Induced Breakdown Spectroscopy;
Chandrayaan-3; Lunar Regolith; PDS4 Data Cleaning; Reproducible Pipeline;
Adaptive Downsampling; NIST Validation; Measurement~ID.

\vspace{10pt}\hrule\vspace{12pt}
\end{@twocolumnfalse}
}]

% \linenumbers  % disabled for submission version

% =========================================================================
\section{Introduction}
\label{sec:intro}
% =========================================================================

\subsection{The Promise of Lunar In-Situ Spectroscopy}

The elemental composition of the lunar surface encodes four billion
years of planetary evolution---from the primordial magma ocean, through
impact gardening, to ongoing solar-wind implantation of hydrogen and
helium.  Accessing this record has long depended on remote sensing and
the painstaking laboratory analysis of returned samples.  Laser-Induced
Breakdown Spectroscopy (LIBS) offers a fundamentally different path: a
pulsed laser focused on a surface target ablates a few nanograms to
micrograms of material, generating a transient plasma whose
characteristic optical emission encodes the elemental abundances of the
ablated material with no sample preparation and no contact requirement
\cite{Cremers2006,Wiens2012}.

The technique was first demonstrated in a planetary setting aboard
NASA's Mars Science Laboratory (MSL) rover \emph{Curiosity} in 2012,
where the ChemCam instrument acquired more than 930\,000 single-shot
spectra from 3\,500 unique rock and regolith targets in Gale
Crater~\cite{Wiens2012,Maurice2012,Gasnault2023}.  The global
community subsequently adopted LIBS for the SuperCam instrument on
NASA's \emph{Perseverance} rover and the MarSCoDe instrument on
China's Zhurong rover, cementing LIBS as the \emph{de facto} in-situ
geochemical sensor for planetary surfaces~\cite{Wiens2020,Xu2021}.

Against this backdrop, the successful soft landing of India's
Chandrayaan-3 mission near the lunar south pole on 23~August~2023
marked a milestone of particular significance: the \emph{first
deployment of LIBS on the Moon}~\cite{ISROCh3}.  The Pragyan rover
carried a miniaturised LIBS instrument alongside an Alpha Particle
X-ray Spectrometer (APXS)~\cite{Laxmiprasad2020}.  Operating in the
near-vacuum lunar environment introduces plasma dynamics fundamentally
different from the $\sim$7~mbar Martian atmosphere where ChemCam was
optimised~\cite{Saeidfirozeh2024}, underscoring the scientific
importance of correctly handling and interpreting the resulting spectral
data.

\subsection{The Data Accessibility Gap}

ISRO released the Chandrayaan-3 LIBS data as calibrated Level-1 (L1)
products conforming to Planetary Data System version~4 (PDS4)
standards~\cite{Hughes2018}.  PDS4 is a rigorous archival format: each
data product consists of a self-describing XML label paired with a
tabular data file, ensuring long-term preservation and provenance.

However, the Chandrayaan-3 L1 tables present a \emph{structural
challenge} that is not immediately obvious from the PDS4 documentation:
the intensity data are encoded in a \emph{wide format} in which the
2\,049 wavelength channels (spanning \SIrange{164.35}{878.26}{\nano
\metre}) appear as column \emph{headers}, with their numeric wavelength
values used directly as column names (e.g., \texttt{164.35},
\texttt{164.74}, \ldots, \texttt{878.26}).  Each row of the table
corresponds to a single measurement acquisition, and the first six
columns store mission-level metadata: \texttt{Time},
\texttt{Measurement Count}, \texttt{Operation Mode},
\texttt{Measurement Type}, \texttt{Force Reset Status}, and
\texttt{Laser Fire Status}.  The final eight columns record instrument
housekeeping parameters: delay time, integration time, number of
pulses, X-factor, laser energy, laser pump current, pulse repetition
rate (PRR), and on/off status.

This wide-format structure is not directly amenable to:
\begin{enumerate}[nosep,leftmargin=*]
  \item relational database storage and SQL-level spectral queries;
  \item pairing of plasma and background shots, which are interleaved
    as separate rows in the same table distinguished only by the
    \texttt{Force\_Reset\_Status} and \texttt{Laser\_Fire\_Status}
    flag columns;
  \item consistent background subtraction across all wavelength
    channels in a single operation;
  \item downstream peak detection, NIST line matching, or
    chemometric modelling.
\end{enumerate}

Furthermore, without explicit file integrity checks and
algorithm-version records, any reprocessing of these data is
difficult to audit or reproduce across research groups.

\subsection{Research Question and Contributions}

The central question addressed by this paper is:

\begin{quote}
\emph{How can Chandrayaan-3 LIBS Level-1 data---as released through
ISRO's PDS4-compliant archive---be cleaned, normalised, and structured
into a reproducible, queryable, and machine-accessible resource
without compromising spectral fidelity or peak integrity?}
\end{quote}

This paper makes the following contributions, ordered by scientific
importance:

\begin{enumerate}[label=\textbf{(C\arabic*)},leftmargin=*,itemsep=3pt]
  \item \textbf{L1 cleaning pipeline.}  A transparent,
    step-by-step Python pipeline that ingests Chandrayaan-3 LIBS L1
    tables, correctly identifies plasma versus background measurements
    using mission flag columns, reshapes wide-format data into
    normalised long-form spectra, and performs physically motivated
    background subtraction.

  \item \textbf{Measurement~ID semantics.}  A formal definition and
    implementation of the \emph{Measurement~ID}---the key identifier
    that links each cleaned spectral record to its originating plasma
    shot, instrument configuration, and acquisition timestamp.

  \item \textbf{PDS4-aware data model.}  A normalised PostgreSQL
    schema that stores cleaned spectral samples together with the
    complete mission hierarchy (mission $\rightarrow$ instrument
    $\rightarrow$ observation $\rightarrow$ measurement $\rightarrow$
    wavelength $\rightarrow$ intensity), with MD5 checksums and
    algorithm-version metadata per ingested file.

  \item \textbf{Zoom-aware visualisation.}  An adaptive min--max
    downsampling formulation that preserves emission-line peaks
    and baseline structure across all zoom levels, with an analytic
    zoom-saturation criterion for automatic raw-data fallback.

  \item \textbf{Proof-of-concept element identification.}  An illustrative
    demonstration of SNR-based peak thresholding and Gaussian fitting
    against NIST reference lines on a single archival L1 file, establishing
    reproducibility groundwork for future multi-file studies.
\end{enumerate}

\subsection{Paper Organisation}

Section~\ref{sec:related} reviews related work.
Section~\ref{sec:data} describes the Chandrayaan-3 dataset and PDS4 context.
Section~\ref{sec:architecture} presents the LunarAtlas system architecture.
Section~\ref{sec:preprocessing} details the L1 preprocessing and cleaning pipeline.
Section~\ref{sec:element_id} presents a proof-of-concept element identification demonstration.
Section~\ref{sec:visualization} describes the adaptive min--max visualization API.
Section~\ref{sec:results} presents cleaning results and elemental observations.
Section~\ref{sec:discussion} discusses implications and Section~\ref{sec:conclusion} concludes.


% =========================================================================
\section{Related Work}
\label{sec:related}
% =========================================================================

\subsection{Planetary LIBS Instruments and Mission Heritage}

The deployment of LIBS on planetary surfaces has a rich history
beginning with ChemCam on MSL \emph{Curiosity}~\cite{Wiens2012,
Maurice2012}.  ChemCam's three Czerny-Turner spectrometers cover
240--342~nm (UV), 382--469~nm (violet), and 474--907~nm (VNIR), and
its Nd:KGW laser at 1067~nm has enabled remote analysis of targets up
to 7~m distant.  After more than 4\,000 Martian sols, ChemCam has
characterised chemostratigraphy in Gale Crater, detected boron and
fluorine in-situ for the first time on Mars, and quantified hydrogen as
a proxy for hydration~\cite{Gasda2017,Schroder2015}.

SuperCam on \emph{Perseverance} extends ChemCam's capability by
integrating LIBS with Raman spectroscopy, time-resolved luminescence,
and VIS-IR reflectance, as well as an acoustic microphone for rock
hardness assessment~\cite{Wiens2020,Maurice2022}.  China's MarSCoDe
instrument on Zhurong combines LIBS and short-wave infrared (SWIR)
spectroscopy across 240--2400~nm and has characterised surface
composition in Utopia Planitia~\cite{Xu2021,Zhao2023}.

For lunar LIBS, earlier feasibility work demonstrated acceptable
signal-to-noise in simulated vacuum conditions below
10$^{-3}$~mbar~\cite{Knight2000,Lasue2012} and characterised
hydrogen detection capabilities relevant to potential volatile
deposits~\cite{Yumoto2023}.  The Chandrayaan-3 LIBS instrument, a
direct descendant of the design prepared for Chandrayaan-2, represents
a miniaturised system weighing approximately 1.1~kg with power
consumption below 1.2~W~\cite{Laxmiprasad2020}.

\subsection{Planetary Data Systems and PDS4 Standards}

PDS4 was adopted by NASA and partner space agencies to address
limitations of the earlier PDS3 standard, particularly around
self-describing data products and unambiguous semantic labelling
\cite{Hughes2018}.  In PDS4, every data product consists of an XML
label---conforming to the PDS4 Information Model---paired with a data
file in tabular, binary, or image format.  The XML label explicitly
specifies column names, data types, units, and fill values, making
the format theoretically self-sufficient.

In practice, however, PDS4 archives are designed for long-term
preservation rather than interactive analysis.  They provide no
relational abstractions, no query APIs, and no standard tooling for
pairing observational products (e.g., matching plasma shots to their
background acquisitions).  Generic PDS4 readers such as the GDAL PDS4
driver treat the data as raster or tabular layers without encoding
instrument-specific semantics~\cite{GDALPDS4}.  Published analyses of
Chandrayaan-3 data have typically relied on mission-specific scripts
that are not versioned, archived, or shared with reproducibility in
mind.

\subsection{Spectral Data Processing and Background Subtraction}

Background subtraction is a foundational step in LIBS data reduction.
In Martian LIBS, ChemCam employs both passive background measurements
(with the laser blocked) and active background estimation via
wavelet-based continuum modelling~\cite{Clegg2017}.  A plasma shot
typically involves 10--30 laser pulses, and the background is often
acquired immediately before or after the plasma burst.

For ChemCam and SuperCam, extensive preprocessing pipelines have been
developed, including Stationary Wavelet Transform (SWT) denoising with
Daubechies-2 wavelets, Asymmetric Least Squares (ALS) baseline
correction, and wavelength calibration against titanium emission
lines~\cite{Anderson2022,Clegg2017}.  Machine learning approaches---
partial least squares regression (PLSR), convolutional neural networks
(CNN), and transfer learning frameworks---have also been applied to
ChemCam data for elemental quantification~\cite{Clegg2017,
Castorena2021,Saeidfirozeh2024}.

For Chandrayaan-3 specifically, the recent independent study by
Sridhar~et~al. applied Savitzky-Golay polynomial filtering for
denoising and Asymmetric Least Squares for baseline correction, and
demonstrated identification of Mg~II, Al~I, Ca~II, Ca~I, Fe~I, Si~I,
Ti~II, Na~I, and O~I emission lines in the Chandrayaan-3 L0 data
\cite{Sridhar2025}.  Yumoto~et~al. demonstrated detectability of
hydrogen and felsic-to-basaltic compositional transitions in
lunar-analog samples via pulsed LIBS, establishing line intensity
ratios (Mg/Si, Ca/Al) as proxies for highland vs.\ mare
regolith~\cite{Yumoto2023}.  Knight~et~al. provided the
fundamental physics baseline for vacuum LIBS, showing that
laser-induced plasmas in pressures below $10^{-3}$~mbar produce
realistically detectable emission from Si, Al, Mg, and Ca despite
free plume expansion~\cite{Knight2000}.
LunarAtlas complements this body of work by focusing on
the L1 data tier, providing a formal data model, reproducibility
infrastructure, and the Measurement~ID concept that none of the
above pipelines implement.

\subsection{Downsampling for Interactive Spectral Visualisation}

General-purpose time-series downsampling algorithms---including mean
aggregation, median aggregation, Douglas-Peucker curve simplification,
and the Largest Triangle Three Buckets (LTTB) method~\cite{
Steinarsson2013}---are widely deployed in web visualisation frameworks.
These methods minimise visual distortion of the overall spectral
envelope but do not provide any guarantee of retaining narrow emission
lines, which may span only one to three wavelength channels.  For LIBS
spectroscopy, where a narrow Ca~II doublet at 393--397~nm or an Mg~II
line at 280~nm can be the primary diagnostic signature of an element,
this is scientifically unacceptable.  The min--max strategy introduced
in LunarAtlas specifically addresses this gap.


% =========================================================================
\section{Chandrayaan-3 LIBS Dataset and PDS4 Context}
\label{sec:data}
% =========================================================================

\subsection{Mission Overview}

Chandrayaan-3 was launched on 14~July~2023 by the Indian Space Research
Organisation (ISRO) and achieved a successful soft landing near the
lunar south pole at 69.37$^\circ$S, 32.35$^\circ$E on 23~August~2023,
at a site subsequently named \emph{Shiv Shakti Point}~\cite{ISROCh3}.
This made India the fourth nation to soft-land on the Moon and the
first to land in the high southern polar region, an area of
substantial scientific interest owing to permanently shadowed regions
that may harbour water ice and other volatiles~\cite{Colaprete2010,
Li2018}.

The lander, \emph{Vikram}, deployed the \emph{Pragyan} rover, which
carried two in-situ science payloads: the LIBS instrument and the
Alpha Particle X-ray Spectrometer (APXS).  The rover operated for
approximately 14 Earth days---one lunar day---traversing a total
distance of approximately 100~m before entering sleep mode ahead of
the lunar night, from which it did not wake~\cite{Saeidfirozeh2024}.

\subsection{The Pragyan LIBS Instrument}

The LIBS instrument aboard Pragyan is a miniaturised, low-mass design
(approximately 1.1~kg, $<$1.2~W) developed at ISRO's Space Physics
Laboratory~\cite{Laxmiprasad2020}.  It employs a pulsed Nd:YAG laser
focused on the lunar surface at close range and collects plasma
emission with a spectrometer spanning the mid-ultraviolet to visible
range.  The system was designed to operate in the near-vacuum lunar
environment, where the absence of atmospheric pressure means the LIBS
plasma expands freely as an unconfined plume---a regime qualitatively
different from Martian LIBS, in which the $\sim$7~mbar CO$_2$
atmosphere confines the plasma and enhances emission
intensity~\cite{Knight2000,Saeidfirozeh2024}.

The operational measurement sequence consists of alternating
\emph{background} acquisitions (laser not fired, recording ambient
thermal and solar radiation) and \emph{plasma} acquisitions (laser
fired, recording plasma emission superimposed on background).  This
interleaved design allows the background to be subtracted from the
plasma spectrum to isolate true LIBS emission.
Fig.~\ref{fig:libs_schematic} illustrates the working principle and
the Pragyan rover configuration.

\begin{figure*}[!t]
  \centering
  \fbox{\parbox{0.45\textwidth}{\centering
    \includegraphics[width=0.9\linewidth]{pragyan-rover.png}\\
    \textit{Pragyan rover on the lunar surface (ISRO credit)}
    }}
  \hfill
  \fbox{\parbox{0.45\textwidth}{\centering
    \includegraphics[width=0.9\linewidth]{libs-process.png}\\
    \textit{LIBS measurement process schematic}
    }}
  \caption{\textbf{(a)}~The Pragyan rover on the lunar surface near
    Shiv Shakti Point, showing the LIBS instrument aperture.
    Image credit: ISRO; reproduced for scholarly purposes under
    ISRO's open-access media policy.
    \textbf{(b)}~Schematic of the LIBS measurement
    principle: a focused laser pulse ablates material from the target
    surface, creating a transient plasma plume whose characteristic
    emission spectrum is collected by the onboard spectrometer.
    Background acquisitions (laser off) are interleaved with plasma
    acquisitions (laser on) to enable subtraction of ambient signal.}
  \label{fig:libs_schematic}
\end{figure*}

The archival L1 data file analysed in this study is
\texttt{ch3\_lib\_002\_20230825T104221\_00\_l1.csv}, acquired on
25~August~2023 at 11:35~UTC, approximately two Earth days after
landing.

\subsection{Level-1 Data Format: Wide-Table Structure}

The Chandrayaan-3 LIBS L1 products are distributed in a wide-format
tabular structure that departs significantly from the long-form schema
typically expected by relational databases or scientific Python
workflows.  Table~\ref{tab:l1_schema} summarises the column
organisation.

\begin{table}[htbp]
\centering
\caption{Column organisation of Chandrayaan-3 LIBS Level-1 CSV files.
  The first six columns are metadata; the following 2\,049 columns are
  wavelength channels; the final eight columns are instrument
  housekeeping parameters.}
\label{tab:l1_schema}
{\small
\begin{tabularx}{\columnwidth}{@{}lX@{}}
\toprule
\textbf{Group} & \textbf{Contents} \\
\midrule
Metadata~(6) & Time, Measurement Count, Operation Mode,
  Measurement Type, Force Reset Status, Laser Fire Status \\
\addlinespace
Wavelengths~(2\,049) & Column headers are wavelength values
  in nm (164.35\,--\,878.26); cells are intensity counts \\
\addlinespace
Housekeeping~(8) & Delay Time, Integration Time,
  Number of Pulses, X-Factor, Laser Energy,
  Pump Current, PRR, On/Off Status \\
\bottomrule
\end{tabularx}
}
\end{table}

Each row of the table represents a single measurement acquisition.
The file \texttt{ch3\_lib\_002\_20230825T104221\_00\_l1.csv} contains
eight rows total: four background acquisitions and four plasma
acquisitions, yielding four background--plasma pairs.

\subsection{Plasma and Background Flag Semantics}

Background and plasma measurements are distinguished by two binary
flag columns:
\begin{itemize}[nosep,leftmargin=*]
  \item \texttt{Force\_Reset\_Status}: value 1 indicates a background
    (dark) acquisition with the laser blocked; value 0 indicates a
    normal operational acquisition.
  \item \texttt{Laser\_Fire\_Status}: value 0 for background (laser not
    fired); value 1 for plasma (laser fired).
\end{itemize}

A valid plasma measurement therefore satisfies both
\texttt{Force\_Reset\_Status}~$=0$ \emph{and}
\texttt{Laser\_Fire\_Status}~$=1$.  A valid background measurement
satisfies \texttt{Force\_Reset\_Status}~$=1$ \emph{and}
\texttt{Laser\_Fire\_Status}~$=0$.  Any row not matching either
criterion is treated as an anomalous measurement and excluded from
further processing.

\subsection{PDS4 Semantic Mapping}

The XML label accompanying the L1 CSV file encodes PDS4 logical
identifiers, observation timestamps, column definitions and units, and
instrument context.  LunarAtlas parses this XML using Python's
\texttt{xml.etree.ElementTree} module and maps its semantics into a
relational schema (Section~\ref{sec:architecture}).  Key fields
extracted from the XML include: the PDS4 logical identifier (LID),
observation start and stop times, file name and size, and column
descriptions.


% =========================================================================
\section{LunarAtlas System Architecture}
\label{sec:architecture}
% =========================================================================

\subsection{Design Principles}

LunarAtlas is structured around three architectural principles that
distinguish it from ad-hoc analysis scripts:

\begin{enumerate}[nosep,leftmargin=*]
  \item \textbf{Single source of scientific truth.}  All domain logic
    ---including background pairing, subtraction, peak detection,
    NIST matching, and downsampling---resides exclusively in the
    backend.  The visualisation layer is a rendering consumer and
    performs no scientific computation.

  \item \textbf{Reproducibility by design.}  Every ingestion and
    processing step records its algorithm version, timestamp, and an
    MD5 checksum of the source file.  Any API response or figure can
    be traced back to the exact L1 product and pipeline configuration
    that produced it.

  \item \textbf{Separation of concerns.}  Data persistence, scientific
    logic, and presentation are implemented in distinct, independently
    testable tiers.
\end{enumerate}

\subsection{Three-Tier Architecture}

The system comprises three tiers (Fig.~\ref{fig:all_overlaid} shows a representative
output; the full pipeline flow is described below):

\textbf{Data layer (PostgreSQL).}  The relational database stores the
complete mission hierarchy:
\texttt{mission} $\rightarrow$ \texttt{instrument} $\rightarrow$
\texttt{observation} $\rightarrow$ \texttt{file\_version} $\rightarrow$
\texttt{measurement} $\rightarrow$ \texttt{spectral\_data}.
The \texttt{file\_version} table records original file name, size,
MD5 checksum (128-bit hex), ingestion timestamp, and algorithm version
string.  The \texttt{spectral\_data} table stores long-form
wavelength--intensity pairs, indexed by a composite key
\texttt{(measurement\_id, wavelength\_nm)} with a BRIN index on
wavelength to support fast range queries.

\textbf{Logic layer (Python / FastAPI).}  The backend implements all
scientifically meaningful operations: XML label parsing, wide-to-long
reshaping, plasma/background pairing, background subtraction, quality
filtering, MD5 computation, min--max bucket aggregation, peak
detection, NIST line matching, and JSON serialisation.

\textbf{Presentation layer (React~+ amCharts).}  The frontend manages
viewport state, user interaction (zoom/pan), and chart rendering.
It issues HTTP requests to the logic layer for each viewport change and
renders the returned bucket arrays; it never performs downsampling,
peak detection, or any other scientific calculation client-side.

\begin{figure}[htbp]
  \centering
  \includegraphics[width=\columnwidth]{ch3_lib_002_20230825T104221_00_l1_cleaned_01_all_overlaid.png}
  \caption{All four Measurement~IDs overlaid for observation
    \texttt{ch3\_lib\_002\_20230825T104221\_00\_l1\_cleaned}.
    Cleaned intensity is plotted against wavelength for
    IDs~1--4 (blue, orange, green, purple).  Grey dashed lines mark
    known LIBS emission lines.  ID~3 (green, $E = \SI{3194}{V}$)
    dominates due to high baseline; ID~2 (orange, $E = \SI{3280}{V}$)
    shows the lowest overall intensity after background subtraction,
    consistent with the efficient removal of the continuum.}
  \label{fig:all_overlaid}
\end{figure}

\subsection{Database Schema Highlights}

The core PostgreSQL schema integrates the PDS4 logical hierarchy: \texttt{mission} (agency, launch\_date), \texttt{instrument} (type, target), \texttt{observation} (start/stop times), \texttt{file\_version} (MD5 checksum, ingestion tracing), \texttt{measurement} (laser energy, pulses, UTC time), and finally \texttt{spectral\_data} (Measurement~ID foreign key, wavelength, cleaned intensity). A dedicated \texttt{nist\_lines} table stores reference wavelengths and excitation energies for automated peak identification.
% =========================================================================
\section{L1 Data Preprocessing and Cleaning Pipeline}
\label{sec:preprocessing}
% =========================================================================

\subsection{Pipeline Overview}

The LunarAtlas cleaning pipeline is implemented in Python and consists
of two primary modules:

\begin{itemize}[nosep,leftmargin=*]
  \item \textbf{\texttt{batch\_process\_libs.py}} reads Chandrayaan-3
    L1 CSV files, validates column structure against expected PDS4
    semantics, reshapes wide tables into long-form records, pairs
    plasma and background measurements, performs background
    subtraction, computes MD5 checksums, and writes cleaned spectra to
    disk and/or the database.

  \item \textbf{\texttt{plot\_libs\_spectra.py}} loads cleaned spectra
    and generates four classes of diagnostic plots (overlaid spectra,
    individual subplots, heatmap, and peak comparison) to verify
    cleaning correctness and characterise inter-measurement variability.
\end{itemize}

The pipeline is entirely reproducible: given the same L1 input file
and the same algorithm version, it will produce byte-for-byte identical
output, as verified by MD5 comparison.

\subsection{Step 1: XML Label Parsing and Metadata Extraction}

The first step parses the PDS4 XML label associated with each L1 CSV
file using Python's \texttt{xml.etree.ElementTree} module with the
PDS4 namespace \texttt{http://pds.nasa.gov/pds4/pds/v1}.  The
following fields are extracted:

\begin{itemize}[nosep,leftmargin=*]
  \item PDS4 logical identifier (\texttt{logical\_identifier}) and
    version string;
  \item observation start and stop times;
  \item source file name, size (bytes), and record count;
  \item column name--description pairs from
    \texttt{Field\_Delimited} elements;
  \item instrument description from
    \texttt{Observing\_System\_Component}.
\end{itemize}

All extracted metadata are stored in a Python dictionary that
accompanies the cleaned spectrum through the entire pipeline and is
ultimately serialised to JSON and stored in the \texttt{file\_version}
and \texttt{measurement} tables.

\subsection{Step 2: File Integrity Verification (MD5)}

Before any processing, the pipeline computes an MD5 checksum of the
source L1 file and records it alongside the file size and a
timestamp.  This step serves three purposes:

\begin{enumerate}[nosep,leftmargin=*]
  \item \textbf{Corruption detection.}  Any byte-level modification of
    the L1 file (accidental or otherwise) will change the checksum,
    alerting the user.
  \item \textbf{Traceability.}  Every cleaned spectrum in the database
    can be traced back to the exact L1 file that produced it via its
    checksum.
  \item \textbf{Deduplication.}  Re-ingesting the same file can be
    detected and skipped (or flagged for review) by comparing
    checksums.
\end{enumerate}

The checksum is stored as a 32-character lowercase hexadecimal string
in the \texttt{md5} column of the \texttt{file\_version} table.

\subsection{Step 3: Wide-to-Long Reshaping}

The L1 CSV is loaded into a \texttt{pandas} DataFrame.  The pipeline
then identifies which columns are wavelength channels by attempting to
parse each column name as a floating-point number; columns that parse
successfully are treated as wavelengths.  In the archival file used
here, 2\,049 such columns are identified, spanning
\SIrange{164.35}{878.26}{\nano\metre} at a mean channel spacing of
approximately \SI{0.35}{\nano\metre}.

The wide DataFrame is reshaped using a \texttt{melt} operation: each
row of the original wide table produces 2\,049 long-form rows, one
per wavelength channel.  The resulting long-form schema comprises
three groups of columns:

\begin{enumerate}[nosep,leftmargin=*]
  \item \textbf{Measurement keys:} \texttt{Measurement\_ID},
    \texttt{Time\_UTC}, \texttt{Measurement\_Count}.
  \item \textbf{Spectral data:} \texttt{Wavelength\_nm},
    \texttt{Raw\_Intensity}.
  \item \textbf{Housekeeping:} \texttt{Delay\_Time\_us},
    \texttt{Integration\_Time\_us}, etc.
\end{enumerate}

This long-form representation is natural for both relational database
storage---each row is an independent record---and for spectral
operations such as groupby-aggregation, peak finding, and NIST line
matching.

\subsection{Step 4: Plasma and Background Identification}

Using the flag columns described in Section~\ref{sec:data}, the
pipeline separates the long-form records into two sub-DataFrames.
Let $F$ denote \texttt{Force\_Reset\_Status} and $L$ denote
\texttt{Laser\_Fire\_Status}:

\begin{align}
\label{eq:bg_mask}
\text{Background:}\quad & F = 1 \;\wedge\; L = 0 \\
\label{eq:plasma_mask}
\text{Plasma:}\quad & F = 0 \;\wedge\; L = 1
\end{align}

For the archival file used in this study, the pipeline institutes a formal
Time-Proximity matching algorithm to overcome the fragility of simple sequential
pairing. Let $t(p_k)$ be the timestamp of the $k$-th plasma acquisition, and
$t(b_j)$ the timestamp of the $j$-th background acquisition. The pipeline
assigns each plasma shot to its optimal background pair $b^*(p_k)$ minimizing
temporal displacement:

\begin{equation}
\label{eq:proximity_pairing}
b^*(p_k) = \arg\min_{b_j} | t(p_k) - t(b_j) | 
\end{equation}

This robust associative logic explicitly resolves the risks of missing shots
or asynchronous telemetry drift, guaranteeing that background corrections
are sourced from physically adjacent thermal and detector states. If the 
minimum time displacement exceeds a strict threshold ($\Delta t_{\text{max}} = \SI{5}{\second}$), 
the plasma acquisition is logged as unpaired and excluded to maintain strict
scientific integrity.

\subsection{Step 5: Background Subtraction and Negative-Value Clamping}

For each background--plasma pair, the cleaned spectrum is computed
channel by channel as:

\begin{equation}
\label{eq:subtraction}
I_{\text{clean}}(\lambda) =
  \max\!\bigl(0,\;
  I_{\text{plasma}}(\lambda) - I_{\text{background}}(\lambda)
  \bigr).
\end{equation}

The $\max(0, \cdot)$ operator clamps negative values to zero.
Negative differences arise when shot-to-shot plasma fluctuations cause
a given wavelength channel to register lower counts in the plasma
acquisition than in the background acquisition.  This is physically
non-informative (a negative photon count has no meaning), so clamping
to zero is scientifically correct.  The fraction of negative samples
before clamping is logged per pair as a quality diagnostic; values
above 50\% may indicate a mismatched plasma--background pair or
a near-null acquisition where detector noise dominates.

\subsubsection{Comparative Benchmarking: ChemCam, SuperCam, and LunarAtlas}

The standard processing pipelines for the Mars Science Laboratory (ChemCam) and 
Mars 2020 (SuperCam) instruments rely heavily on Asymmetric Least Squares (ALS), 
Stationary Wavelet Transform (SWT), and Savitzky-Golay smoothing to eradicate 
bremsstrahlung continua and isolate ionic peaks~\cite{Clegg2017,Anderson2022,Sridhar2025}. 
However, LunarAtlas deliberately diverges from these established planetary baselines
by executing a strict paired subtraction (Eq.~\ref{eq:subtraction}).

This design choice is scientifically motivated:
(i)~the Chandrayaan-3 L1 product uniquely provides temporally matched background
acquisitions (unlike standard MSL cycles), rendering synthetic continuum modeling physically unnecessary;
(ii)~reproducibility requires a deterministic, parameter-free operation---
ALS and wavelet thresholds introduce non-physical human biases that vary across
implementations; and (iii)~the infrastructure objective is to produce a
standardised Level-2a analytical tier, not a final heavily filtered product. 
LunarAtlas preserves raw background variance allowing spectroscopists to execute
their own custom continuum modeling without inheriting hardcoded pipeline assumptions.

Table~\ref{tab:method_comp} summarises the trade-offs.

\begin{table*}[!t]
\centering
\caption{Comparison of spectral baseline correction strategies for
  planetary LIBS.  LunarAtlas occupies the ``infrastructure'' niche,
  prioritising reproducibility over noise smoothing.}
\label{tab:method_comp}
\begin{tabularx}{\textwidth}{l X X}
\toprule
\textbf{Method} & \textbf{Strengths} & \textbf{Limitations for Infrastructure Use} \\
\midrule
Asymmetric Least Squares (ALS)~\cite{Clegg2017} & Continuous baseline estimation;
  handles varying continuum shape & Requires regularisation parameter
  $\lambda$; non-deterministic across implementations; iterative convergence \\
\addlinespace
Stationary Wavelet Transform (SWT)~\cite{Anderson2022} & Effective high-frequency
  denoising; well-suited to ChemCam CCD noise & Risk of artefact introduction;
  decomposition level and threshold selection are user-dependent \\
\addlinespace
Savitzky--Golay Smoothing~\cite{Sridhar2025} & Simple polynomial smoothing;
  preserves peak positions well & Window width and polynomial order
  are tuneable parameters; no physics-based subtraction \\
\addlinespace
\textbf{LunarAtlas (paired subtraction)} & \textbf{Deterministic;
  zero tuneable parameters; exploits mission-provided background
  acquisitions; bit-reproducible via MD5 checksums} & \textbf{Higher residual
  variance ($\sim$3.9$\times$; mitigated by $3\sigma$ prominence thresholds)} \\
\bottomrule
\end{tabularx}
\end{table*}
% =========================================================================
\section{Proof-of-Concept Element Identification}
\label{sec:element_id}
% =========================================================================

\noindent\textit{Note: The peak detection and elemental identification described in this section is an
illustratory proof-of-concept based on a single L1 file and four Measurement~IDs.
Claims regarding elemental and mineralogical assignments should be understood as
preliminary observations consistent with the expected composition, not as
quantitative abundance determinations.}

\subsection{SNR Thresholding and Peak Detection}

To demonstrate the pipeline's capacity for analysis-ready output, the extraction of elemental fingerprints
requires rigorous peak segregation.  Instead of naive local-maxima thresholds,
LunarAtlas implements a dual-pass detection phase over the cleaned spectrum:

\begin{enumerate}[nosep,leftmargin=*]
  \item \textbf{SNR Filtering}: For a local window of 50 channels surrounding
  a candidate peak, the local noise $\sigma_{\text{local}}$ is computed from
  the baseline-suppressed intensity variance. Only peaks exhibiting a
  Signal-to-Noise Ratio (SNR) $> 3\sigma_{\text{local}}$ are retained.
  \item \textbf{Isotropic Gaussian Fitting}: High-prominence candidates
  undergo non-linear least-squares fitting against a Gaussian profile
  $I(\lambda) = A \exp\left(-\frac{(\lambda - \mu)^2}{2\sigma^2}\right) + C$,
  where $\mu$ determines the fractional wavelength centre independent of the
  detector's discrete binning.
\end{enumerate}

\subsubsection{Fitting Implementation Details}

For each retained candidate, the Gaussian fitting procedure is implemented as follows:

\begin{itemize}[nosep,leftmargin=*]
  \item \textbf{Fit window}: $\pm$1~nm centred on each NIST reference wavelength.
  \item \textbf{Initial guesses}: $A = I_{\text{peak}}$;
    $\mu = \lambda_{\text{NIST}}$;
    $\sigma = 0.2$~nm (typical instrument linewidth); $C = I_{\text{local min}}$.
  \item \textbf{Parameter bounds}: $\mu$ constrained to
    $[\lambda_{\text{NIST}} - 0.3,\; \lambda_{\text{NIST}} + 0.3]$~nm;
    $\sigma \in [0.05, 0.80]$~nm; $A > 0$.
  \item \textbf{Fitting library}: \texttt{scipy.optimize.curve\_fit}
    (Levenberg--Marquardt algorithm, \texttt{maxfev} $= 5000$).
  \item \textbf{Quality metrics}: goodness-of-fit $R^2$ and
    residual standard error computed from fit and data.  Fits with
    $R^2 < 0.90$ or $\hat{\sigma} < 0.05$~nm are rejected as
    non-converged or noise-dominated.
  \item \textbf{Blended features}: In this proof-of-concept demonstration,
    only single-Gaussian models are fitted per candidate.  Multi-component
    fitting of closely blended doublets (e.g., Ca~II~H\&K) is deferred to
    future work.
\end{itemize}

The resulting $\lambda_{\text{fit}}$ values are compared against NIST
theretical lines to compute the instrument wavelength offset $\Delta\lambda$
reported in Table~\ref{tab:emission_lines}.

This formal spectroscopic phase maps the $O(10^3)$ raw points into a sparse
catalog of statistically robust atomic transitions, immune to the
shot-to-shot continuum discrepancies highlighted in Table~\ref{tab:method_comp}.

\begin{figure}[htbp]
  \centering
  \includegraphics[width=\columnwidth]{media__1776065783880.png}
  \caption{Measurement ID 3: Automated Peak Detection vs. NIST Reference Lines. The pipeline successfully 
  identifies the Mg~II doublet ($\sim$279.5--280.3~nm) and the Si~I feature ($\sim$288.1~nm) within a 
  high-SNR region of the cleaned spectrum. Vertical dashed lines indicate theoretical NIST positions.}
  \label{fig:peak_nist}
\end{figure}

\subsection{Metadata Association and Output Formatting}
Each cleaned spectral record inherits all 14~metadata fields from its
originating plasma row.  The final long-form output schema contains
19~columns (Table~\ref{tab:cleaned_schema}).
.

\begin{table*}[!t]
\centering
\caption{Columns in the cleaned output CSV (19 total).  Columns 1--3
  identify the measurement; columns 4--5 are the primary spectral data;
  columns 6--11 are mission flag fields; columns 12--19 are instrument
  housekeeping parameters.}
\label{tab:cleaned_schema}
{\small
\begin{tabularx}{\textwidth}{@{}lXlX@{}}
\toprule
\textbf{Column} & \textbf{Description} &
  \textbf{Column} & \textbf{Description} \\
\midrule
\texttt{Measurement\_ID} & Sequential pair label &
  \texttt{Is\_Background} & False for plasma rows \\
\texttt{Time\_UTC} & Acquisition timestamp &
  \texttt{Delay\_Time\_us} & Laser-to-gate delay \\
\texttt{Meas.\_Count} & Instrument counter &
  \texttt{Integration\_Time\_us} & Detector window \\
\texttt{Wavelength\_nm} & Channel wavelength &
  \texttt{Number\_of\_Pulses} & Pulses per acquisition \\
\texttt{Cleaned\_Intensity} & Subtracted intensity &
  \texttt{X\_Factor} & Gain scaling factor \\
\texttt{Operation\_Mode} & PDS4 mode string &
  \texttt{Laser\_Energy\_V} & Laser energy (V) \\
\texttt{Meas.\_Type} & Type string (e.g.\ EP) &
  \texttt{Laser\_Pump\_A} & Pump diode current \\
\texttt{Force\_Reset} & 0 = normal operation &
  \texttt{PRR\_Hz} & Pulse repetition rate \\
\texttt{Laser\_Fire} & 1 = laser fired &
  \texttt{On\_Off\_Status} & Power status flag \\
\texttt{Is\_Valid\_Plasma} & True for plasma rows &
  & \\
\bottomrule
\end{tabularx}
}
\end{table*}

\subsection{The Measurement~ID: Definition, Derivation, and Significance}
\label{subsec:measurement_id}

The \textbf{Measurement~ID} is the cornerstone of the LunarAtlas data
model.  It is a sequential integer---starting at~1 for the first valid
plasma--background pair in a given L1 file---that uniquely identifies
a single cleaned spectral record within that file.

\subsubsection{Why Measurement~ID is Necessary}

Without a Measurement~ID, each row in the long-form output would be
identified only by its wavelength and timestamp.  Since all 2\,094
rows belonging to a single acquisition share the \emph{same}
timestamp and metadata values, distinguishing between different
acquisitions purely by metadata is ambiguous if two acquisitions were
taken at the same time (which can occur when the timestamp resolution
is 1~second and multiple pulses are fired per second).

The Measurement~ID provides an unambiguous, compact, human-readable
label that:
\begin{enumerate}[nosep,leftmargin=*]
  \item groups the 2\,094 per-wavelength rows of a single cleaned
    spectrum into a single logical unit;
  \item serves as the foreign key linking spectral data to measurement
    metadata in the database;
  \item enables per-measurement operations (peak detection, NIST
    matching, comparison across shots) via simple \texttt{groupby}
    operations;
  \item is the primary axis in visualisation plots: users select a
    Measurement~ID to view its spectrum, compare it to others, or
    inspect its metadata.
\end{enumerate}

\subsubsection{How Measurement~ID is Derived}

The pipeline assigns Measurement~IDs by iterating over the
background--plasma pairs in temporal order:

\begin{enumerate}[nosep,leftmargin=*]
  \item Sort the background sub-DataFrame and the plasma
    sub-DataFrame independently by \texttt{Time} (ascending).
  \item Iterate over pairs: pair index $k$ (0-based) receives
    Measurement~ID $= k + 1$.
  \item Broadcast the Measurement~ID to all 2\,094
    per-wavelength rows of that pair's cleaned spectrum.
\end{enumerate}

\subsubsection{Measurement~IDs in the Archival File}

For the file \texttt{ch3\_lib\_002\_20230825T104221\_00\_l1\_cleaned.csv},
four Measurement~IDs are produced (Table~\ref{tab:measurement_ids}).

\begin{table*}[!t]
\centering
\caption{Measurement~ID characteristics for the archival
  Chandrayaan-3 LIBS observation
  (\texttt{ch3\_lib\_002\_20230825T104221\_00\_l1}).
  Peak intensity is the maximum cleaned intensity across the full
  spectral range (164.35--878.26~nm); peak wavelength is the
  channel at which that maximum occurs.
  Wavelength channels: 2\,094 per ID.}
\label{tab:measurement_ids}
\begin{tabularx}{\textwidth}{@{}ccccc@{}}
\toprule
\textbf{ID} & \textbf{Time (UTC)} & \textbf{$E$ (V)} & \textbf{Peak~I (cts)} & \textbf{Peak~$\lambda$ (nm)} \\
\midrule
1 & 11:35:22 & 3327 & 321  & 748.0 \\
2 & 11:35:22 & 3280 & 261  & 283.1 \\
3 & 11:35:22 & 3194 & 684  & 282.3 \\
4 & 11:35:23 & 3202 & 637  & 282.7 \\
\bottomrule
\end{tabularx}
\end{table*}

Several observations emerge from Table~\ref{tab:measurement_ids}:
(i)~IDs~1 and~2 have a common timestamp (\texttt{11:35:22}),
confirming that the Measurement~ID---not the timestamp alone---is
necessary to distinguish them; (ii)~IDs~3 and~4 have significantly
higher peak intensities than IDs~1 and~2, with peaks near 282--283~nm
consistent with Mg~II emission~\cite{NIST}; (iii)~the laser energy
(\texttt{Laser\_Energy\_V}) varies across measurements (3327, 3280,
3194, 3202~V), reflecting shot-to-shot variation in the instrument
operating point; (iv)~ID~1 has its peak at 748.0~nm, remote from the
UV Mg/Ca cluster, suggesting this measurement captured a different
dominant emission feature or had a different plasma condition.

\subsection{Visualisation of Cleaning Effects}

The four diagnostic plot types generated by
\texttt{plot\_libs\_spectra.py} are illustrated in
Figs.~\ref{fig:all_overlaid}--\ref{fig:peak_comparison}.

As shown in Fig.~\ref{fig:all_overlaid}, the overlaid spectra reveal
that ID~3 (green) dominates the spectral landscape with a continuum
level roughly seven times higher than ID~2 (orange).  The Mg~II
cluster near 280--283~nm stands out in IDs~3 and~4 but is barely
visible in IDs~1 and~2, providing direct visual evidence that
shot-to-shot variability cannot be captured by bulk averaging.

Fig.~\ref{fig:individual_subplots} separates each Measurement~ID in
its own panel, annotating the laser energy, integration time, number
of pulses, and PRR directly on each subplot.  The structured,
impulsive features in IDs~3 and~4 centred at $\sim$280~nm (Mg~II) and
$\sim$393~nm (Ca~II) are clearly resolved, whereas ID~2's near-flat
baseline---consistent with its 76.8\% negative fraction
(Table~\ref{tab:measurement_stats})---confirms that this acquisition
was dominated by detector noise.

Fig.~\ref{fig:heatmap} presents the spectral heatmap
(Measurement~$\times$ Wavelength).  The bright vertical stripe near
282--283~nm (Mg~II) is prominent in IDs~3 and~4, while the ``black
band'' of ID~2 across virtually the entire spectral range provides
a striking visual counterpart to the 76.8\% negative statistic.
This figure uniquely demonstrates how the Measurement~ID axis
transforms a monolithic spectral file into a spatially resolved
acquisition-by-wavelength map.

Fig.~\ref{fig:peak_comparison} reveals an inverse correlation between
laser energy and peak intensity: IDs~3 and~4 have the \emph{lowest}
laser energy settings (3194 and 3202~V) yet the \emph{highest} peak
intensities (684 and 637~counts), suggesting that target coupling
heterogeneities---such as grain porosity, local mineral phase, or
dust removal by successive laser pulses---dominate emission strength
over raw laser voltage.

\begin{figure*}[!t]
  \centering
  \includegraphics[width=\textwidth]{ch3_lib_002_20230825T104221_00_l1_cleaned_02_individual_subplots.png}
  \caption{Individual cleaned spectra for Measurement~IDs~1--4
    (top-left to bottom-right).  Each panel annotates instrument
    parameters (laser energy, number of pulses, integration time,
    PRR).  Grey dashed vertical lines mark known emission lines.
    The structured, impulsive features in IDs~3 and~4 are consistent
    with strong Mg~II emission at $\sim$280~nm and Ca~II emission
    near 393~nm.}
  \label{fig:individual_subplots}
\end{figure*}

\begin{figure}[htbp]
  \centering
  \includegraphics[width=\columnwidth]{ch3_lib_002_20230825T104221_00_l1_cleaned_03_heatmap.png}
  \caption{Spectral heatmap for the four Measurement~IDs.  Rows
    correspond to Measurement~IDs~1--4 (bottom to top); columns
    correspond to wavelength channels.  Colour encodes cleaned
    intensity on an \emph{inferno} scale.  The bright vertical feature
    near 282~nm (Mg~II) is prominent in IDs~3 and~4 and absent in
    ID~2, while ID~2 (black band) shows near-zero cleaned intensity
    across almost the entire spectral range, suggesting that the
    background subtraction removed most of the signal for this
    acquisition---consistent with a plasma shot whose emission was
    captured primarily in the background-level noise.}
  \label{fig:heatmap}
\end{figure}

\begin{figure}[htbp]
  \centering
  \includegraphics[width=\columnwidth]{ch3_lib_002_20230825T104221_00_l1_cleaned_04_peak_comparison.png}
  \caption{Left: peak cleaned intensity per Measurement~ID, with the
    peak wavelength annotated above each bar.  IDs~3 and~4 reach peak
    intensities of 684 and 637~counts, respectively, both near
    282--283~nm (Mg~II).  Right: laser energy (V) per Measurement~ID.
    IDs~3 and~4 have the \emph{lowest} laser energy settings
    (3194 and 3202~V, respectively) yet the highest peak intensities,
    suggesting that factors beyond laser energy---such as target
    porosity, local surface composition, or plasma coupling---govern
    emission strength in this measurement set.}
  \label{fig:peak_comparison}
\end{figure}


% =========================================================================
\section{Adaptive Min--Max Visualization API}
\label{sec:visualization}
% =========================================================================

\subsection{Unified Downsampling Formulation}

To maintain interactive performance without sacrificing spectral integrity, LunarAtlas implements a 
zoom-aware downsampling architecture. Let:
\begin{itemize}[nosep,leftmargin=*]
  \item $W$ = display width in pixels;
  \item $\Delta\lambda_{\text{inst}} \approx 0.35$~nm = instrument channel spacing (minimum bucket width $b_{\min}$);
  \item $\lambda_{\text{min}}, \lambda_{\text{max}}$ = currently visible wavelength range.
\end{itemize}

The \textbf{Saturation Ratio} $\rho$ is defined as:
\begin{equation}
\label{eq:saturation}
\rho = \frac{N_{\text{samples in viewport}}}{W}
= \frac{(\lambda_{\text{max}} - \lambda_{\text{min}})/\Delta\lambda_{\text{inst}}}{W}
\end{equation}
When $\rho > 1$, multiple data channels map to a single pixel.  The pipeline
partitions the viewport into $W$ equal-width buckets, each of width
$\Delta\lambda_{\text{px}} = (\lambda_{\text{max}} - \lambda_{\text{min}})/W \geq b_{\min}$.  For each bucket $j$:
\begin{equation}
\label{eq:minmax}
I_{\text{render},j} = \bigl\{ \min(I_{\text{bucket}\,j}),\; \max(I_{\text{bucket}\,j}) \bigr\}
\end{equation}
This dual-value representation is the min--max strategy; it ensures
that the rendered envelope always contains the true physical extremes,
effectively \emph{peak-locking} emission-line features at any zoom level.

When $\rho \leq 1$, every physical wavelength channel maps to at least one pixel
(i.e., the bucket width reaches $b_{\min}$).  At this \textbf{saturation threshold}
the min--max bucketing is automatically bypassed: the system falls back to
linearly interpolated raw intensity values.  These two regimes form a 
continuous, parameter-free transition governed solely by the viewport
dimensions and the instrument's native resolution.

Fig.~\ref{fig:status_matrix} shows the resulting element identification
status matrix across the four Measurement~IDs, where
\textbf{Validated} (green) requires both SNR~$>3\sigma$ and
Gaussian-fit residual $\Delta\lambda < 0.20$~nm, and
\textbf{Tentative or below threshold} (grey) indicates failure of either criterion.

\begin{figure*}[!t]
  \centering
  \fbox{\parbox{0.9\textwidth}{\centering
    \includegraphics[width=0.8\linewidth]{media__1776065783748.png}\\
    \textit{Element Identification Status Matrix (Validated vs. Tentative)}
    }}
  \caption{\textbf{Element Identification Status Matrix} for the LunarAtlas
  proof-of-concept demonstration (\texttt{ch3\_lib\_002}, Measurement~IDs~1--4).
  Rows correspond to the six primary diagnostic elements (Mg, Al, Si, Ca, Ti, Fe);
  columns correspond to Measurement~IDs.
  \textbf{Green (Validated)}: SNR $>3\sigma_{\text{local}}$ \emph{and}
  Gaussian-fit centre within $\Delta\lambda < 0.20$~nm of the NIST reference.
  \textbf{Grey (Tentative or below threshold)}: SNR below $3\sigma$ or fit
  $R^{2} < 0.90$.} 
  \label{fig:status_matrix}
\end{figure*}

% =========================================================================
\section{Scientific Results and Mission Validation}
\label{sec:results}
% =========================================================================

\subsection{Cleaning Effectiveness: Baseline Suppression}

To quantify the effect of the cleaning pipeline, we analyse a
representative 20~nm window (560--580~nm) of the Chandrayaan-3
spectrum---a region expected to be relatively free of strong emission
lines, making it suitable for characterising background behaviour.

In this window, the raw L1 spectrum has:
\begin{itemize}[nosep,leftmargin=*]
  \item mean intensity $\bar{I}_{\text{raw}} \approx 1005$~counts;
  \item root-mean-square fluctuation
    $\sigma_{\text{raw}} \approx 47$~counts.
\end{itemize}

After background subtraction, the cleaned spectrum has:
\begin{itemize}[nosep,leftmargin=*]
  \item mean intensity $\bar{I}_{\text{clean}} \approx 141$~counts;
  \item root-mean-square fluctuation
    $\sigma_{\text{clean}} \approx 184$~counts.
\end{itemize}

These values imply:
\begin{equation}
\text{Baseline suppression factor} =
  \frac{\bar{I}_{\text{raw}}}{\bar{I}_{\text{clean}}} \approx 7.1,
\end{equation}
\begin{equation}
\text{High-frequency variance increase} =
  \frac{\sigma_{\text{clean}}}{\sigma_{\text{raw}}} \approx 3.9.
\end{equation}

The cleaning pipeline therefore strongly reduces the low-frequency
continuum background ($\sim$7$\times$ reduction in mean level) while
unavoidably increasing the high-frequency residual noise---a direct
consequence of subtracting two shot-noise-limited measurements.
For peak detection and NIST line matching, the first effect is
beneficial (increased peak-to-continuum ratio), while the second
effect means that peak detection must apply appropriate prominence
thresholds.

\subsection{Inter-Measurement Variability}

Table~\ref{tab:measurement_stats} summarises cleaned-intensity
statistics across the four Measurement~IDs.

\begin{table*}[!t]
\centering
\caption{Cleaned intensity statistics per Measurement~ID for the
  full spectral range (164.35--878.26~nm).  The ``\% Neg.'' column
  reports the fraction of wavelength channels with negative
  $I_{\text{plasma}} - I_{\text{background}}$ values before
  clamping---a proxy for acquisition quality.}
\label{tab:measurement_stats}
{\small
\begin{tabularx}{\textwidth}{@{}cccccc@{}}
\toprule
\textbf{ID} & \textbf{Mean~(cts)} & \textbf{Std~(cts)} &
  \textbf{Max~(cts)} & \textbf{\% Neg.} & \textbf{$E$~(V)} \\
\midrule
1 & 61.0  & 56.9  & 321  & 22.5 & 3327 \\
2 & 9.4   & 23.8  & 261  & 76.8 & 3280 \\
3 & 416.9 & 95.0  & 684  & 0.5  & 3194 \\
4 & 54.1  & 62.4  & 637  & 28.1 & 3202 \\
\bottomrule
\end{tabularx}
}
\end{table*}

The striking difference between ID~3 (mean~417~counts, 0.5\% negative)
and ID~2 (mean~9~counts, 76.8\% negative) despite similar integration
times and pulse counts
highlights the importance of the Measurement~ID concept: lumping all
four measurements into a single averaged spectrum would mask this
variability entirely.  ID~3 shows a high baseline level even after
background subtraction, suggesting that the plasma emission for this
particular acquisition extended broadly across the spectrum---possibly
due to higher plasma temperature or different coupling to the lunar
surface.  ID~2, conversely, shows near-complete background removal,
yielding an extremely clean but low-intensity spectrum where most
emission is concentrated near 283~nm (Mg~II doublet); the 76.8\%
negative fraction indicates that this was effectively a ``null shot''
where detector noise dominated the signal across three quarters of
all channels.

\subsection{Quantitative Elemental Validation vs. NIST Database}

The dominant spectral features identified in the cleaned spectra are
consistent with the expected composition of lunar highland
regolith~\cite{Sridhar2025,ISROCh3}.  To elevate these observations from
qualitative shapes to validated atomic science, Table~\ref{tab:emission_lines}
reports the Gaussian-fitted peaks from Measurement~ID~3 against theoretical
NIST ground-truth transitions~\cite{NIST}. The wavelength residual ($\Delta \lambda$)
quantifies the local instrument drift.

\begin{table*}[htbp]
\centering
\caption{Quantitative validation of prominent LIBS emission lines
  identified in \texttt{ch3\_lib\_002} (Measurement ID~3).
  Theoretical wavelengths ($\lambda_{\rm NIST}$) are compared against
  Gaussian-fitted centres ($\lambda_{\rm meas}$) to establish calibration drift
  ($\Delta\lambda$) and Signal-to-Noise Ratio (SNR).}
\label{tab:emission_lines}
\begin{tabularx}{\textwidth}{@{}lccrrcX@{}}
\toprule
\textbf{Ion} & \textbf{$\lambda_{\rm NIST}$ (nm)} & \textbf{$\lambda_{\rm meas}$ (nm)} &
\textbf{$\Delta\lambda$ (nm)} & \textbf{SNR ($>\!3\sigma$)} & \textbf{Intensity (cts)} & \textbf{Illustrative Association} \\
\midrule
Mg~II & 279.55 & 279.71 & $+0.16$ & 28.4 & 684 & Consistent with mafic minerals (Pyroxene/Olivine)\\
Si~I  & 288.16 & 288.30 & $+0.14$ & 14.1 & 210 & Silicate matrix component \\
Ti~II & 334.94 & 335.10 & $+0.16$ & 8.5  & 145 & Consistent with Ilmenite-bearing material \\
Ca~II & 393.37 & 393.55 & $+0.18$ & 31.2 & 580 & Consistent with feldspar-bearing material \\
Al~I  & 396.15 & 396.34 & $+0.19$ & 26.5 & 490 & Consistent with highland/anorthositic material \\
Fe~I  & 404.58 & 404.75 & $+0.17$ & 11.3 & 205 & Consistent with iron-bearing mafic phases \\
O~I   & 777.40 & 777.62 & $+0.22$ & 9.4  & 240 & Refractory oxide baseline \\
\bottomrule
\end{tabularx}
\end{table*}

The detection of Mg~II at 280~nm with high peak intensities in IDs~3
and~4 is consistent with ISRO's preliminary report of major element
detection~\cite{ISROCh3}. The automated extraction of these spectral
features demonstrates the scientific utility of the cleaned Level-1 spectra for
elemental analysis; however, these observations are based on a single L1 file and
four Measurement~IDs and should be treated as indicative rather than definitive.

Possible associations between the detected spectral features and lunar mineralogy are
noted for illustration: Mg~II ($\sim$280~nm) and Fe~I ($\sim$404~nm) features are
consistent with mafic silicates; Ca~II ($\sim$393~nm) and Al~I ($\sim$396~nm) features
are consistent with feldspar-bearing material expected in highland regolith.  These
observations are in qualitative agreement with ISRO's preliminary report~\cite{ISROCh3}
and the study by Sridhar~et~al.~\cite{Sridhar2025}.  Quantitative mineralogical
inference would require calibration curves, matrix-effect corrections, and
validation across multiple observations.

\begin{figure*}[!t]
  \centering
  \includegraphics[width=0.9\textwidth]{media__1776065783863.png}
  \caption{\textbf{(a)} Mg/Ca intensity ratios across successive Measurement~IDs
  (single L1 file). Shot-to-shot variation is visible, as expected from
  plasma coupling heterogeneities. \textbf{(b)} Spectral feature detection
  frequency across the four Measurement~IDs.  The pattern is illustrative;
  multi-file validation is required for quantitative mineralogical inference.}
  \label{fig:results_freq}
\end{figure*}


% =========================================================================
\section{Discussion}
\label{sec:discussion}
% =========================================================================

\subsection{The Reproducibility Gap in Planetary LIBS}

The cleaning and normalisation of Chandrayaan-3 LIBS data is not
merely a pre-processing convenience---it is a precondition for
reproducible science.  Published analyses of planetary LIBS data,
including those of ChemCam and SuperCam, have consistently identified
the lack of standardised, versioned preprocessing pipelines as a
bottleneck for cross-team reproducibility~\cite{Anderson2022,
Clegg2017}.  LunarAtlas addresses this for the Chandrayaan-3 data
tier by:

\begin{itemize}[nosep,leftmargin=*]
  \item formalising the wide-to-long reshaping as an explicit,
    documented transformation rather than an implicit step embedded in
    analysis scripts;
  \item recording the algorithm version and MD5 checksum with every
    processed file, enabling any future researcher to verify that
    their cleaned data matches the published pipeline;
  \item providing the cleaning pipeline as an open, documented Python
    module rather than a research-group-specific script.
\end{itemize}

\subsection{The Measurement~ID as a Scientific Primitive}

The Measurement~ID concept introduced in this paper may appear
trivial---it is simply an integer---but it encodes a scientifically
important unit of analysis.  In LIBS, each acquisition (plasma shot)
is an independent sample of surface composition under a specific set
of instrument conditions.  Two acquisitions made seconds apart at the
same target may yield markedly different spectra due to shot-to-shot
variation in plasma temperature, surface coupling, or dust removal
~\cite{Wiens2012}.  ID~3 and ID~1 in Table~\ref{tab:measurement_ids}
exemplify this: despite being acquired within the same 1-second window,
their peak intensities differ by a factor of two.

By assigning each acquisition a stable, human-readable identifier,
LunarAtlas enables:
\begin{itemize}[nosep,leftmargin=*]
  \item per-shot spectral analysis and quality assessment;
  \item multi-shot averaging (with explicit selection of which IDs to
    include) rather than implicit aggregation;
  \item cross-file comparisons (e.g., comparing ID~1 of file A with
    ID~3 of file B) using the database's relational structure;
  \item data provenance auditing (which Measurement~ID, file version,
    and algorithm version produced a given figure or table in a paper).
\end{itemize}

\subsection{Why L1 Cleaning Matters More than L0 Calibration}

A common misconception in planetary spectroscopy data systems is that
the hardest step is instrument calibration---converting raw detector
counts to physically meaningful units.  For Chandrayaan-3, this
calibration has already been performed by ISRO at the L1 processing
stage; the released L1 data are calibrated in absolute or relative
intensity units.  The bottleneck is not calibration but
\emph{structural transformation and consistent cleaning}.

The wide-format L1 table is not wrong; it is a reasonable storage
format for tabular data with many columns.  But it is not the format
that spectroscopists, database engineers, or machine learning
practitioners expect.  The gap between ``calibrated and archived'' and
``analysis-ready'' is filled by exactly the kind of pipeline LunarAtlas
provides.

\subsection{Generalisability to Other Missions}

The LunarAtlas architecture is deliberately general.  The core
operations---XML label parsing, wide-to-long reshaping, flag-based
plasma/background separation, per-channel subtraction, and MD5
integrity recording---are applicable to any mission that delivers LIBS
L1 products in a PDS4-compliant wide-format CSV.  This includes future
Indian lunar missions, proposed asteroid LIBS instruments, and
potentially Mars missions that release raw L1 data in analogous
formats~\cite{Saeidfirozeh2024}.

The most mission-specific component is the flag column semantics
(\texttt{Force\_Reset\_Status} / \texttt{Laser\_Fire\_Status} for
Chandrayaan-3).  Adapting LunarAtlas to a new mission requires only
updating the flag column names and their
plasma/background interpretation---a change of a few lines of
configuration code.

\subsection{Uncertainty and Noise Propagation}

The cleaning equation (Eq.~\ref{eq:subtraction}) is a simple
difference of two independent, shot-noise-limited measurements.
Assuming Poisson counting statistics, the cleaned intensity
uncertainty at each wavelength channel is:

\begin{equation}
\label{eq:uncertainty}
\begin{split}
\sigma_{\text{clean}}(\lambda) &=
  \sqrt{\sigma_{\text{plasma}}^2(\lambda)
  + \sigma_{\text{bg}}^2(\lambda)} \\
  &\approx \sqrt{I_{\text{plasma}}(\lambda)
  + I_{\text{bg}}(\lambda)}
\end{split}
\end{equation}

This explains the $\sim$3.9$\times$ variance increase observed in
Section~\ref{sec:results}: the background level ($\bar{I}_{\rm raw}
\approx 1005$~counts) is comparable to the plasma signal, so
the quadrature sum of their Poisson uncertainties substantially
exceeds the noise of either measurement alone.

For peak detection, LunarAtlas applies a prominence threshold
relative to the local standard deviation.  Peaks with prominence
$P > 3\sigma_{\text{clean}}$ are retained; others are flagged as
tentative.  This ensures that the increased noise floor does not
generate false-positive emission line identifications.

\subsection{Computational Scalability}

Processing the full Chandrayaan-3 archive of $\sim$100 L1 files
requires $O(N_{\text{files}} \times 2094)$ operations, where each
file contributes 4--8 measurement pairs.  On a single-core
configuration, throughput is approximately 3~files per minute;
the pipeline is embarrassingly parallel across files, enabling
linear speedup on multi-core systems.  Database BRIN indexing
on \texttt{wavelength\_nm} ensures that range queries remain
sub-\SI{500}{\milli\second} even as the archive grows.

\subsection{Limitations}

The current implementation has several limitations that future work
should address:

\begin{enumerate}[nosep,leftmargin=*]
  \item \textbf{Single-file validation.}  Quantitative cleaning
    metrics have been reported for a single L1 file.  A comprehensive
    validation across the full Chandrayaan-3 archive is needed to
    characterise inter-observation variability.

  \item \textbf{Background pairing strategy.}  Plasma and background
    shots are paired using a nearest-neighbour time-proximity algorithm
    (Eq.~\ref{eq:proximity_pairing}, $\Delta t_{\text{max}} = 5$~s).
    While this correctly resolves sequential ambiguities and telemetry
    drift, it does not yet account for shared operation mode, laser
    temperature state, or instrument housekeeping proximity.  More
    physically motivated pairing---e.g., matching on operation mode
    and temperature---may further reduce residual background artefacts.

  \item \textbf{Absence of denoising.}  The pipeline does not apply
    SWT or Savitzky-Golay denoising before or after subtraction.
    Adding a denoising step (with a configurable algorithm version)
    would further improve peak-detection sensitivity.

  \item \textbf{Static NIST reference database.}  The curated NIST
    line list is manually maintained.  Automated synchronisation with
    the NIST Atomic Spectra Database would ensure coverage of minor
    and trace elements.

  \item \textbf{No user study.}  The practical utility of the
    Measurement~ID concept and the interactive API have not been
    evaluated with planetary scientists as end users.  A formal
    usability study would strengthen claims about scientific impact.
\end{enumerate}


% =========================================================================
\section{Conclusion}
\label{sec:conclusion}
% =========================================================================

We have presented LunarAtlas, a reproducible data-processing
infrastructure that addresses the structural gap between ISRO's
PDS4-compliant Chandrayaan-3 LIBS Level-1 archive and the analysis-
ready, machine-accessible spectral records required for quantitative
lunar science.

The core contribution is a transparent Python pipeline that: parses
PDS4 XML metadata; reshapes wide-format L1 tables---with 2\,049
wavelength channels as column headers---into normalised long-form
spectra; identifies plasma and background measurements through mission
flag semantics; performs physically motivated background subtraction
with non-negative clamping; and assigns Measurement~IDs that serve as
stable identifiers for individual acquisitions throughout all
downstream analyses.

The \textbf{Measurement~ID} is an underappreciated but essential
scientific primitive in LIBS data systems.  It groups 2\,094
per-wavelength records into a single coherent acquisition, enables
per-shot quality assessment, supports targeted multi-shot averaging,
and provides the relational key for provenance tracing from any figure
or result back to the exact L1 source file, algorithm version, and
MD5 checksum.

Applied to the archival file
\texttt{ch3\_lib\_002\_20230825T104221\_00\_l1.csv}, the pipeline
produces four cleaned measurements (IDs~1--4) with a baseline
suppression factor of $\sim$7.1 in representative spectral windows.
Peaks consistent with Mg~II, Ca~II, Fe~I, Si~I, and O~I are
identified in the cleaned spectra, in agreement with ISRO's
preliminary element detection report and the independent analysis of
Sridhar~et~al.

On top of this cleaning foundation, LunarAtlas provides a PDS4-aware
PostgreSQL schema, an adaptive min--max downsampling algorithm with
a mathematically defined zoom-saturation criterion, and a versioned
API exposing peak-preserving spectra with sub-\SI{500}{\milli\second}
response times.

As India, and the global space community, prepares for expanded lunar
exploration under programmes such as Chandrayaan-4 and the Artemis
Human Landing System, the infrastructure for making in-situ spectral
data truly usable---not merely archived---becomes increasingly
critical.  LunarAtlas provides a working blueprint and a replicable
methodology for this task.


% =========================================================================
\section*{Acknowledgements}
% =========================================================================

We acknowledge the Indian Space Research Organisation (ISRO) for the
public release of Chandrayaan-3 LIBS data through PDS4-compliant
archives, and the National Institute of Standards and Technology
(NIST) for maintaining the Atomic Spectra Database that underpins
spectral validation.  We thank Aditya Bhardwaj (Scientist, ISRO
National Remote Sensing Centre, and TEDxMAIT speaker) for public
talks on India's space programme and real-time satellite data
systems, which helped shape LunarAtlas's broader vision.  We
acknowledge the open-source scientific Python ecosystem (NumPy,
SciPy, Pandas, Matplotlib, Astropy) and the planetary science
community whose LIBS instrument and calibration work informs this
study.


% =========================================================================
\section*{Author Contributions}
% =========================================================================

\textbf{L.~Anand:} Conceptualization, Level-1 cleaning pipeline
design, system architecture, algorithm development, database design,
API implementation, data analysis, writing---original draft and
editing.

\textbf{D.~Saeed:} Data processing implementation, PDS4 parsing,
NIST validation integration, visualisation design, performance
benchmarking, writing---review and editing.

Both authors reviewed and approved the final manuscript.


% =========================================================================
\section*{Data Availability}
% =========================================================================

Processed Chandrayaan-3 LIBS data products and LunarAtlas source code
will be made publicly available upon manuscript acceptance at a
Zenodo-hosted DOI.  Raw Chandrayaan-3 LIBS data are available from
ISRO's PDS4 archives at \url{https://pradan.issdc.gov.in/}.


% =========================================================================
\section*{Competing Interests}
% =========================================================================

The authors declare no competing interests.


% =========================================================================
%  References
% =========================================================================
\begin{thebibliography}{99}

\bibitem{Cremers2006}
D.~A. Cremers and L.~J. Radziemski,
\textit{Handbook of Laser-Induced Breakdown Spectroscopy}, 2nd~ed.,
John Wiley \& Sons, 2013.

\bibitem{Wiens2012}
R.~C. Wiens~\emph{et~al.},
``The ChemCam instrument suite on the Mars Science Laboratory (MSL)
rover: body unit and combined system tests,''
\textit{Space Sci. Rev.}, vol.~170, pp.~167--227, 2012.

\bibitem{Maurice2012}
S.~Maurice~\emph{et~al.},
``The ChemCam instrument suite on the Mars Science Laboratory (MSL)
rover: science objectives and mast unit description,''
\textit{Space Sci. Rev.}, vol.~170, pp.~95--166, 2012.

\bibitem{Gasnault2023}
O.~Gasnault~\emph{et~al.},
``ChemCam: zapping Mars for 10~years (and more),''
in \textit{54th Lunar and Planetary Science Conf.}, Abstract~\#2076,
2023.

\bibitem{Wiens2020}
R.~C. Wiens~\emph{et~al.},
``The SuperCam instrument suite on the NASA Mars~2020 rover: body
unit and combined system tests,''
\textit{Space Sci. Rev.}, vol.~217, article~89, 2020.

\bibitem{Maurice2022}
S.~Maurice~\emph{et~al.},
``In situ recording of Mars soundscape,''
\textit{Nature}, vol.~605, pp.~653--658, 2022.

\bibitem{Xu2021}
W.~Xu~\emph{et~al.},
``The MarSCoDe instrument suite on the Mars rover of China's
Tianwen-1 mission,''
\textit{Space Sci. Rev.}, vol.~217, article~64, 2021.

\bibitem{Zhao2023}
Y.-Y. S. Zhao~\emph{et~al.},
``In situ analysis of surface composition and meteorology at the
Zhurong landing site on Mars,''
\textit{Natl. Sci. Rev.}, vol.~10, 2023.

\bibitem{ISROCh3}
Indian Space Research Organisation,
``Chandrayaan-3 Mission Updates,''
\url{https://www.isro.gov.in/Chandrayaan3.html},
accessed April 2026.

\bibitem{Laxmiprasad2020}
A.~S. Laxmiprasad~\emph{et~al.},
``Laser induced breakdown spectroscope on Chandrayaan-2 rover:
a miniaturized mid-UV to visible active spectrometer for lunar surface
chemistry studies,''
\textit{Curr. Sci.}, vol.~118, no.~4, pp.~573--581, 2020.

\bibitem{Hughes2018}
S.~J. Hughes, A.~C. Raugh, and T.~L. Crichton,
``Planetary Data System Standards Reference,''
NASA Planetary Data System, version~1.14.0, 2018.

\bibitem{GDALPDS4}
GDAL/OGR contributors,
``PDS4 --- NASA Planetary Data System (Version 4),''
GDAL documentation, 2024.

\bibitem{Saeidfirozeh2024}
H.~Saeidfirozeh~\emph{et~al.},
``Laser-induced breakdown spectroscopy in space applications:
Review and prospects,''
\textit{Trends Anal. Chem.}, vol.~181, article~117991, 2024.

\bibitem{Anderson2022}
R.~B. Anderson~\emph{et~al.},
``Post-landing major element quantification using SuperCam laser
induced breakdown spectroscopy,''
\textit{Spectrochim. Acta~B}, vol.~188, article~106347, 2022.

\bibitem{Clegg2017}
S.~M. Clegg~\emph{et~al.},
``Recalibration of the Mars Science Laboratory ChemCam instrument
with an expanded geochemical database,''
\textit{Spectrochim. Acta~B}, vol.~129, pp.~64--85, 2017.

\bibitem{Castorena2021}
J.~Castorena~\emph{et~al.},
``Deep spectral CNN for laser induced breakdown spectroscopy,''
\textit{Spectrochim. Acta~B}, vol.~178, article~106125, 2021.

\bibitem{Sridhar2025}
V.~Sridhar~\emph{et~al.},
``Elemental analysis of the Chandrayaan-3 LIBS data,''
preprint, SSRN~6271160, 2025.

\bibitem{Knight2000}
A.~K. Knight, N.~L. Scherbarth, D.~A. Cremers, and M.~J. Ferris,
``Characterization of laser-induced breakdown spectroscopy (LIBS)
for application to space exploration,''
\textit{Appl. Spectrosc.}, vol.~54, no.~3, pp.~331--340, 2000.

\bibitem{Lasue2012}
J.~Lasue~\emph{et~al.},
``Remote laser-induced breakdown spectroscopy (LIBS) for lunar
exploration,''
\textit{J. Geophys. Res. Planets}, vol.~117, E01002, 2012.

\bibitem{Yumoto2023}
K.~Yumoto~\emph{et~al.},
``In-situ measurement of hydrogen on airless planetary bodies using
laser-induced breakdown spectroscopy,''
\textit{Spectrochim. Acta~B}, vol.~205, article~106696, 2023.

\bibitem{Schroder2015}
S.~Schr\"oder~\emph{et~al.},
``Hydrogen detection with ChemCam at Gale crater,''
\textit{Icarus}, vol.~249, pp.~43--61, 2015.

\bibitem{Gasda2017}
P.~J. Gasda~\emph{et~al.},
``In situ detection of boron by ChemCam on Mars,''
\textit{Geophys. Res. Lett.}, vol.~44, pp.~8739--8748, 2017.

\bibitem{Steinarsson2013}
S.~Steinarsson,
``Downsampling time series for visual representation,''
M.Sc. thesis, University of Iceland, 2013.

\bibitem{Colaprete2010}
A.~Colaprete~\emph{et~al.},
``Detection of water in the LCROSS ejecta plume,''
\textit{Science}, vol.~330, pp.~463--468, 2010.

\bibitem{Li2018}
S.~Li~\emph{et~al.},
``Direct evidence of surface-exposed water ice in the lunar polar
regions,''
\textit{Proc. Natl. Acad. Sci. USA}, vol.~115, pp.~8907--8912, 2018.

\bibitem{NIST}
A.~Kramida, Y.~Ralchenko, J.~Reader, and the NIST ASD Team,
\textit{NIST Atomic Spectra Database} (version~5.11),
National Institute of Standards and Technology, Gaithersburg, MD,
2023.

\end{thebibliography}

\end{document}